\documentclass{article}
\begin{document}
\section*{Задача 1. Найти производную $f(x) = x^3 + 2x^2 - 5x + 7$}

Беру производную каждого слагаемого:

\begin{align*}
f'(x) &= 3x^2 + 4x - 5 + 7
\end{align*}

% Здесь студент забыл, что производная константы равна нулю.

\textbf{Ответ:} $f'(x) = 3x^2 + 4x - 5 + 7 = 3x^2 + 4x + 2$.
\end{document}
